%==============================================================================
% CHAPTER 2
%==============================================================================
\chapter{Architecture of the AutoOFFBEAT agent}
\markboth{Chapter 2: Architecture of the AutoOFFBEAT agent}{}

\section{Overview}

The architecture adopted is that of a single supervisor orchestrating
specialised tools (figure~\ref{fig:archi}). The user expresses a request in
natural language in a web interface. The supervisor, built on a language model,
selects and chains the necessary tools. Each tool performs a concrete action and
returns a textual report that the model interprets.

\begin{figure}[htbp]
\centering
\begin{tikzpicture}[
  node distance=0.55cm and 0.7cm,
  box/.style={draw,rounded corners=2pt,align=center,font=\small,
              minimum height=0.95cm,inner sep=4pt},
  % font=\footnotesize: "offbeat\_executor" in \small exceeds 2.9 cm and
  % overlaps the neighbouring box.
  tool/.style={box,fill=blue!5,draw=cClad!70,text width=2.9cm,
               font=\footnotesize},
  kb/.style={box,fill=orange!8,draw=cFuel!70,text width=2.7cm,font=\scriptsize},
  sup/.style={box,fill=gray!10,draw=black!60,text width=4.6cm},
  ext/.style={box,fill=none,draw=black!45,dashed,text width=2.9cm},
  fl/.style={-{Latex[length=2mm]},draw=black!65},
  ]

  \node[box,text width=3.6cm,fill=green!5,draw=green!45!black] (user)
    {User\\{\scriptsize natural-language request}};
  \node[sup,below=0.7cm of user] (sup)
    {\textbf{LLM supervisor}\\{\scriptsize selects and chains the tools}};

  \node[tool,below=1.5cm of sup,xshift=-4.9cm] (t1)
    {\texttt{input\_creator}\\{\scriptsize builds the case}};
  \node[tool,below=1.5cm of sup,xshift=-1.65cm] (t2)
    {\texttt{offbeat\_executor}\\{\scriptsize meshes, runs,\\ \textbf{repairs}}};
  \node[tool,below=1.5cm of sup,xshift=1.65cm] (t3)
    {\texttt{data\_processor}\\{\scriptsize post-processes}};
  \node[tool,below=1.5cm of sup,xshift=4.9cm] (t4)
    {\texttt{safety\_analyzer}\\{\scriptsize safety margins}};

  \node[ext,below=1.35cm of t2] (solver)
    {OFFBEAT solver\\{\scriptsize OpenFOAM}};
  \node[ext,below=1.35cm of t3] (vtk)
    {VTK fields\\{\scriptsize $T$, $\sigma$, gap}};

  \node[kb,left=0.55cm of solver,xshift=-0.4cm] (errkb)
    {\texttt{error\_kb.json}\\{\scriptsize crash patterns}};
  \node[kb,right=0.55cm of vtk,xshift=0.4cm] (safekb)
    {\texttt{safety\_kb.json}\\{\scriptsize safety thresholds}};

  \draw[fl] (user) -- (sup);
  \draw[fl] (sup.south) -- ++(0,-0.35) -| (t1.north);
  \draw[fl] (sup.south) -- ++(0,-0.35) -| (t2.north);
  \draw[fl] (sup.south) -- ++(0,-0.35) -| (t3.north);
  \draw[fl] (sup.south) -- ++(0,-0.35) -| (t4.north);
  \draw[fl] (t2) -- (solver);
  \draw[fl] (t3) -- (vtk);
  \draw[fl] (errkb) -- (t2);
  \draw[fl] (safekb) -- (t4);
  \draw[fl,dashed] (solver.east) -- (vtk.west);
\end{tikzpicture}
\caption[AutoOFFBEAT architecture]{Architecture of AutoOFFBEAT.}
\label{fig:archi}
\end{figure}

Three principles underpin this architecture.

\textbf{The language model does not compute.} Every physical quantity comes
from the solver or from \textbf{deterministic} post-processing. This separation
guarantees that the validity of the results does not depend on the reliability
of the language model, which could otherwise hallucinate values.

\textbf{Deterministic first, model as fallback.} Wherever a diagnosis can be
obtained from an explicit rule (an error pattern, a quantified threshold), that
rule takes precedence. The language model only intervenes on failure, and its
opinion is never applied automatically.

\textbf{Domain knowledge lives outside the code.} Error patterns and safety
criteria are stored in editable JSON files. A domain user can enrich the base
without touching the code or rebuilding anything.

\section{The language model factory}

So as not to tie the project to a particular provider, access to the language
model goes through a single factory driven by a configuration file. Three
providers are supported: fully local execution through Ollama, and two remote
services. The rest of the code is unaware of which one is in use.

Two things followed from this choice. The whole development could
be carried out \textbf{at no inference cost}, with a seven-billion-parameter
model running locally. Second, the factory distinguishes the supervisor model
from the \textit{debugging} model: it becomes possible to entrust routine
routing to a small, fast model and crash analysis to a more capable one.

\section{The domain tools}

\subsection{Case generation}

Generating a complete OFFBEAT case from scratch is not realistic. The solver's
main dictionary alone runs to several hundred lines of coupled models, and a
language model would inevitably produce invalid combinations. The strategy
adopted is therefore that of the \textbf{template}: the tool copies a validated
reference case, then substitutes only the requested parameters.

A mapping table associates understandable names (linear heat rate, pellet
radius, simulated duration) with the actual locations in the case files. This
method is essential: it gives the language model a stable and restricted
vocabulary, without requiring it to know the syntax of OpenFOAM dictionaries.
This matters all the more because a local model with few parameters is used,
which can quickly make mistakes. The complete list of exposed parameters is
given in appendix~\ref{ann:cas}.

Two templates are currently available: a one-dimensional PWR rod, used for all
the results in this report, and a two-dimensional axisymmetric $(r,z)$ rod taken
from the OFFBEAT verification cases. The latter was validated end to end
(creation, parameter application, meshing, execution and safety analysis), which
confirmed that the mesh-region-based processing does not depend on the
dimensionality of the case.

\subsection{Solver execution}

The execution tool chains mesh generation and then the solver run, relying on a
maximum-duration safeguard so that a diverging simulation cannot block the
system indefinitely. It first checks that the target directory really is an
OpenFOAM case, and analyses the run log afterwards.

\subsection{Post-processing}

Post-processing reads the solver outputs in VTK format and extracts the useful
quantities: axial and radial temperature profiles, hoop stress, maximum cladding
temperature. It also produces figures. One subtlety, which turned out to be a
significant source of errors (§\ref{sec:bugs}), lies in the organisation of the
data: the complete mesh contains both the pellet and the cladding, and an
extremum computed over the whole domain does not necessarily say anything about
the region of interest.

\section{Self-healing}
\label{sec:selfhealing}

This is the agent's most distinctive contribution. When the solver crashes, the
execution tool attempts to diagnose the cause and remedy it. It implements a
two-stage strategy, shown in figure~\ref{fig:selfheal}.

\begin{figure}[htbp]
\centering
\begin{tikzpicture}[
  node distance=0.5cm,
  b/.style={draw,rounded corners=2pt,align=center,font=\small,
            minimum height=0.85cm,inner sep=4pt,text width=3.4cm},
  d/.style={draw,diamond,aspect=2.4,align=center,font=\scriptsize,
            inner sep=1pt,text width=2.3cm},
  fl/.style={-{Latex[length=2mm]},draw=black!65,font=\scriptsize},
  ]
  \node[b,fill=blue!5] (run) {Run the solver};
  \node[d,below=of run,fill=gray!8] (ok) {success?};
  \node[b,below=1.0cm of ok,fill=green!8,draw=green!45!black] (done)
    {Results ready};
  \node[d,right=1.5cm of ok,fill=gray!8] (match) {known pattern?};
  \node[b,above=1.0cm of match,fill=orange!8,draw=cFuel] (fix)
    {Apply the correction\\{\scriptsize knowledge base}};
  \node[b,right=1.6cm of match,fill=red!5,draw=cAlert] (llm)
    {Diagnosis by the model\\{\scriptsize proposed, not applied}};
  \node[b,below=1.0cm of llm,fill=gray!8] (stop)
    {Stop: user\\ intervention required};

  \draw[fl] (run) -- (ok);
  \draw[fl] (ok) -- node[left]{yes} (done);
  \draw[fl] (ok) -- node[above]{no} (match);
  \draw[fl] (match) -- node[right]{yes} (fix);
  \draw[fl] (fix.north) -- ++(0,0.35) -| node[pos=0.25,above]
    {retry ($\leq 3$)} (run.north);
  \draw[fl] (match) -- node[above]{no} (llm);
  \draw[fl] (llm) -- (stop);
\end{tikzpicture}
\caption[Self-healing loop]{Self-healing loop. Deterministic pattern-based
diagnosis is attempted first; the language model only intervenes as a last
resort and its diagnosis is never applied automatically.}
\label{fig:selfheal}
\end{figure}

The first stage matches the run log against a pattern base. Each entry
associates a regular expression with a plain-language diagnosis and a scripted
correction (reduce the time step, clean the written time directories to force a
clean restart, adjust the simulated duration). This mechanism is fast, free and
reproducible.

The second stage only intervenes if no pattern matches: the end of the log is
then submitted to the language model for interpretation. Crucially, this
diagnosis is \textbf{presented to the user but never applied automatically}: a
human must have the final say on error-sensitive actions.

The number of attempts is capped at three. Chapter~4 shows that this
architecture succeeds on the numerical errors it was designed to handle, but
reaches its limits when faced with \textit{physical} errors. This is one of the
main lessons of the internship.

\section{The documentary assistant}

A final tool, optional and similar to the one in AutoFLUKA, makes the OFFBEAT
documentation available to the agent through semantic search
(\textit{retrieval-augmented generation}). Documents are split into chunks,
embedded by a local embedding model, and indexed. The agent can thereby answer
documentation questions by quoting real excerpts rather than producing plausible
but unverified statements. This tool is designed to disable itself silently if
unavailable: it must never prevent basic operation.

\clearpage
